The simulation study was set up as a host-side validation environment for the network-layer logic. The main
goal was to evaluate forwarding behavior under controlled channel and mobility conditions, without getting
confused by embedded timing artifacts. We did not create a separate protocol just for testing. Instead, we
kept the same network-layer decision process used in the real deployment and ran it in a virtualized
communication setup. This way, any behavior we observed came from the routing and suppression policies, not
from having different code paths for production vs testing.

The evaluation covered both channel/transport-level effects and network-layer forwarding semantics. The
simulation modeled virtual time progression, topology changes, and configurable channel impairments, while the
forwarding pipeline at each node stayed unchanged. So the traces can be read as direct evidence
of how duplicate handling, relay delay policy, and hop-limited dissemination behave in reproducible scenarios.
We followed a staged approach: first basic delivery correctness, then forwarding-policy validation, and then
determinism and scalability stress.

Time evolution was controlled using an explicit virtual clock, so host scheduling jitter did not affect causal
ordering. The scenario only moved forward when we deliberately advanced time. This let us test not only
"eventual delivery" but also verify that delivery does not happen too early when insufficient virtual time has
elapsed. Event processing followed a stable ordering policy, so concurrent transmissions, deferred relays, and
queued channel actions were resolved consistently across repeated runs. This helped separate protocol behavior
from platform variability and supported strict regression comparisons.

The scenarios allowed controlled configuration of channel timing, impairment intensity, congestion behavior,
and randomization seed, along with per-node initial state and mobility evolution. Mobility was represented as
kinematic progression with scheduled state updates, so we could repeat transitions between disconnected and
connected topologies. Since topology updates and event processing were integrated into a single ordered step
sequence, we could also assert the exact virtual time when connectivity changed.

The channel model combined distance-dependent attenuation with receiver sensitivity gating and optional
stochastic impairments. Reception depended on whether the signal conditions were admissible at the receiver.
On top of that, optional packet-error and congestion-driven loss mechanisms introduced controlled non-ideal
behavior. For the stochastic parts, pseudo-random decisions were seed-controlled and tied to event identity.
So with a fixed seed we got identical outcomes, and with different seeds we got intentional divergence. This
balance let us do repeatable verification and sensitivity studies without modifying the forwarding logic.

Validation was done using a reusable harness that captured application-visible deliveries at each node. We
checked both payload content and the overall dissemination structure. Assertions were made using stepped
virtual-time progression, which allowed precise checks on temporal causality, deferred forwarding, and
timeout-related behavior. In parallel, we collected aggregate channel and delivery metrics to quantify
retransmission activity, collision effects, packet-error losses, congestion drops, latency trends, and
utilization. Using both trace-level and aggregate evidence increased confidence in functional correctness and
system-level behavior.

Baseline results confirmed expected fan-out for broadcast dissemination, preserved payload integrity, and
correctly excluded sender self-delivery. Reachability followed the configured radio conditions: nearby
receivers were served when the signal budget was admissible, while distant receivers were excluded until
mobility brought them into effective communication range. This showed that basic delivery behavior matched
geometric and sensitivity constraints before considering higher-level forwarding policies.

Temporal and channel-access evaluations showed that delivery depends on advancing virtual time, and that
contention leads to deferred transmissions rather than uncontrolled immediate rebroadcasts. With concurrent
access, overlapping activity produced measurable collision-related failures, while controlled deferral and
retry behavior still kept progression bounded. Propagation-delay experiments also showed physically consistent
ordering where nearer receivers can observe arrival before farther receivers when the transmission origin is
the same.

Impairment and load experiments verified that packet-error and congestion mechanisms affected delivery in the
expected directions. Strict error thresholds and forced congestion reduced successful receptions and increased
drop-related counters, while light-load regimes below congestion thresholds preserved delivery continuity.
So the simulation can represent both nominal and stressed channel conditions in a way that is easy to
interpret.

Most importantly, simulation-based validation of forwarding semantics confirmed correct behavior for
lifetime-bounded dissemination, directional relay eligibility, multi-hop reach extension, and duplicate
suppression at the application-delivery boundary. Lifetime constraints prevented stale forwarding beyond the
configured validity window. Directional policy constrained relay participation based on propagation intent.
Relaying extended delivery to nodes outside direct source reach. When redundant copies arrived through
different paths, delivery stayed idempotent at the receiver interface, showing effective suppression of
duplicate application exposure.

Scalability and repeatability tests provided additional evidence of robustness. Large-node broadcast workloads
maintained the expected dissemination coverage, and higher-node stepping workloads stayed within practical
runtime budgets for host-based evaluation. Repeating runs with identical seeds produced stable dissemination
snapshots, while changing the seed produced controlled stochastic differences. So we could use the framework
for deterministic regression testing, and also for controlled variability analysis, within one methodology.

The results should be interpreted within the model assumptions. Propagation was represented using analytic
attenuation with configurable impairment terms rather than a fully site-calibrated empirical urban channel
model. Also, contention was abstracted at the event level instead of modeling complete physical-layer waveform
interaction. So, the study provides strong evidence of network-layer logical correctness, timing consistency,
and resilience under modeled conditions, but it does not claim full physical-fidelity prediction for every
deployment environment.

So, the simulation campaign gives coherent evidence of logical correctness, timing consistency, and
forwarding behavior under modeled nominal and adverse conditions. It also establishes a reproducible baseline
that is complemented by the real-world measurements in the next subsection.